\documentclass{beamer}
\usepackage{amssymb}
\usepackage{amsmath}
\usepackage{commath}
\usepackage{amsthm}
\usepackage{lmodern}
\usepackage{scrextend}
\usepackage{listings}
\usepackage{xcolor}
\newcommand*\diff{\mathop{}\!\mathrm{d}}
\usetheme{Luebeck}
\useoutertheme[subsection=false, footline=authortitle]{miniframes}
\usecolortheme{crane}

\lstdefinestyle{base}{
  language=Python,
  emptylines=1,
  breaklines=true,
  basicstyle=\ttfamily\color{black},
  moredelim=**[is][\color{orange}]{@}{@},
}

\title{Visualising weak convergence for mathematical models in shallow geophysical domains}
\subtitle{A problem}
\author{N\'ora Juh\'asz}

\date{18th February, 2019}

\begin{document}

\begin{frame}
\titlepage
\end{frame}

\section{State of the research project}

\begin{frame}{The main ideas of the model}
  \begin{itemize}
    \item Geophysical fluid $\to$ shallow domain characteristics --- the domain is greater horizontally than in the vertical dimension.
    \item Classical / general concepts as foundation:
    \begin{itemize}
      \item The wind flow being modelled by the incompressible Navier-Stokes equations
      \item Using a convection - diffusion equation to describe the pollution effect
    \end{itemize}
  \end{itemize}
\end{frame}

%% coordinate system - axes are independent from the wind.
%% density is one,
%% different viscosities depending H/V
\begin{frame}{The governing equations}
  \begin{itemize}
    \item Fluid velocity equations --- incompressible Navier-Stokes equations with the density taken to be identically equal to one.
      
      $$\partial_t \mathbf{v} + (\mathbf{v} \cdot \nabla)\mathbf{v} -\Delta_{\nu} \mathbf{v} + \nabla q + 2 \mathbf{w} \times \mathbf{v}= \mathbf{g} \text{ in } \Omega_{\epsilon} \times (0,T)$$
    
      $$\nabla \cdot \mathbf{v} = 0 \text{ in } \Omega_{\epsilon} \times (0,T)$$

    \item A convection-diffusion equation representing pollutants in the atmosphere:
    
      $$\partial_t C + \mathbf{v} \cdot \nabla C = K_x \partial_{xx} ^2 C + K_y  \partial_{yy} ^2  C + K_z  \partial_{zz} ^2  C + S$$

  \end{itemize}
\end{frame}

\begin{frame}{Scaling}
%% largely different horizontal and vertical scales
%% epsilon incorporates the shallowness of the domain
%% viscosity rescaling - deduced heuristically from the dimensions
One of the main tools to exploit the "almost" 2D structure of the atmosphere $\to$ introduce the aspect ratio $$\epsilon = \frac{\text{characteristic depth}}{\text{characteristic width}},$$ which incorporates the shallowness in a mathematical way. Natural scaling process:

  \begin{enumerate}
    \item $x=x_1, y=x_2, z=\epsilon x_3$
    \item $v_1 = u_1^{\epsilon}, v_2 = u_2^{\epsilon}, v_3 = \epsilon u_3^{\epsilon}, p = p^{\epsilon},$
    \item $\nu_x = \nu_1, \nu_y = \nu_2 , \nu_z = \epsilon ^ 2\nu_3,$
    \item $K_x = K_1, K_y = K_2, K_z = \epsilon ^ 2 K_3.$
  \end{enumerate}
\end{frame}

\begin{frame}{Local domain and boundary conditions}
  We use a local cartesian coordinate system so that we can neglect the curvature of the Earth.
  \begin{figure}
    \includegraphics[width=\linewidth]{domain.jpg}
    \caption{A local slice of the atmosphere}
    \label{fig:atmosphere}
  \end{figure}
\end{frame}

\begin{frame}{Anisotropic equations}

  In $\Omega \times (0,T)$ we have the following system.

    \[ \partial_t u^\epsilon _1 + \mathbf{u}^\epsilon \cdot \nabla u^\epsilon _1 -\Delta_\nu u^\epsilon _1 -\alpha u_2^\epsilon + \epsilon \beta u_3^\epsilon + \partial_1 p^\epsilon = 0 \]
    
    \[ \partial_t u_2^\epsilon + \mathbf{u}^\epsilon \cdot \nabla u_2^\epsilon -\Delta_\nu u_2^\epsilon + \alpha u_1^\epsilon + \partial_2 p^\epsilon = 0 \]
    
    \[\epsilon^2 \{  \partial_t u_3^\epsilon + \mathbf{u}^\epsilon \cdot \nabla u_3^\epsilon -\Delta_\nu u_3^{\epsilon} \} - \epsilon \beta u_1^\epsilon + \partial_3 p^\epsilon = 0 \]
    
    \[ \nabla \cdot \mathbf{u}^\epsilon= 0  \]
    
    \[ \partial_t C^\epsilon + \mathbf{u}^\epsilon \cdot \nabla C^\epsilon = K_1 \frac{\partial ^2 C^\epsilon}{\partial x_1 ^ 2} + K_2 \frac{\partial ^2 C^\epsilon}{\partial x_2 ^ 2} + K_3 \frac{\partial ^2 C^\epsilon}{\partial x_3 ^ 2} + S_\epsilon \]
    
\end{frame}

%% note that the limit model is 3D - this is a different approach compared to the shallow water depth-integrated model which is 2D
%% in the vertical direction we have the hydrostatic pressure instead of the wind momentum equation
\begin{frame}{Anisotropic equations - Part 2}

  The initial and boundary conditions are chosen to be

    \[ \mathbf{u}^\epsilon = 0 \text{ on } \Gamma \times (0,T), \]
    
    \[ \mathbf{u}^\epsilon (\cdot, t = 0) = \mathbf{u}_0, C^\epsilon (\cdot, t=0) = C^\epsilon_0  \text{ in } \Omega, \]

    $$C^\epsilon = 0 \text{ on } ( \Gamma_U \cup \Gamma_L) \times (0,T), \partial_2 C^\epsilon = \partial_3 C^\epsilon = 0 \text{ on } \Gamma_G \times (0,T),$$
    
    $$\partial_1 C^\epsilon \in L^2 L^2 (0,T; \Gamma_G). $$

\end{frame}

\begin{frame}{Hydrostatic equations}
  
      \[ \partial_t u_1 + \mathbf{u} \cdot \nabla u_1 - \Delta_\nu u_1 -\alpha u_2 + \partial_1 p = 0 \text{ in } \Omega \times (0,T) \]
    
    \[ \partial_t u_2 + \mathbf{u} \cdot \nabla u_2 - \Delta_\nu u_2 -\alpha u_1 + \partial_2 p = 0 \text{ in } \Omega \times (0,T) \]
    
    \[ \partial_3 p = 0 \text{ in } \Omega \times (0,T) \]
    
    \[ \nabla \cdot \mathbf{u} = 0 \text{ in } \Omega \times (0,T) \]
    
    \[ \partial_t C + \mathbf{u} \cdot \nabla C = K_1 \frac{\partial ^1 C}{\partial x_1 ^ 2} + K_2 \frac{\partial ^2 C}{\partial x_2 ^ 2} + K_3 \frac{\partial ^2 C}{\partial x_3 ^ 2} + S \text{ in } \Omega \times (0,T)\]

\end{frame}
    
\begin{frame}{Hydrostatic equations - Part 2}
    \[ u_1 = u_2 = u_3 n_3 = 0 \text{ on } \Gamma \times (0,T) \]
    
    \[ u_i (\cdot, t=0) = u_{0 i}, C (\cdot, t=0) = C_0 \text{ in } \Omega, \text{$i$ = 1,2} \]
    
    $$C = 0 \text{ on } ( \Gamma_U \cup \Gamma_L) \times (0,T), \partial_2 C = \partial_3 C = 0 \text{ on } \Gamma_G \times (0,T),$$
    
    $$ \partial_1 C \in L^2 L^2 (0,T; \Gamma_G) $$
  
\end{frame}

\begin{frame}{Main theorem}

  We show a rigorous connection between the two sets --- taking the limit in the solutions we arrive to a justification of the hydrostatic model for the case of the polluted atmosphere.

  \begin{theorem}

  Let $\mathbf{u}_0$ $\in$ $L^2 (\Omega)^3$, with $\nabla \cdot \mathbf{u}_0 = 0, \mathbf{u}_0 \cdot \mathbf{n} = 0$ on $\partial \Omega$; and let $C(0)\in$ $L^2 (\Omega)^3$, $C(0) = 0$ on $ \partial \Omega$; then as the aspect ratio $\epsilon$ tends to zero, any weak solution $(\mathbf{u}^\epsilon, C^\epsilon)$ of the anisotropic equations converges to a weak solution ($\mathbf{u}$, C) of the hydrostatic equations of the polluted atmosphere.

  \end{theorem}

\end{frame}

\begin{frame}{Energy inequality}
    \changefontsizes{10pt}
  \begin{equation*}
    \begin{split}  
    & \frac{1}{2} (\norm{\mathbf{u}_H^\epsilon (t)}_{L^2}^2 + \epsilon ^ 2 \norm{u_3^\epsilon (t)}_{L^2}^2 + \norm{C^\epsilon (t)}_{L^2}^2) + \int_0 ^ t \big [\norm{\nabla_\nu \mathbf{u}_H^\epsilon}_{L^2}^2 + \epsilon ^ 2 \norm{\nabla_\nu u_3^\epsilon}_{L^2}^2 \big ] \diff \tau \\
    & + \int_0 ^ t \bigg[ (\epsilon K_1 - \zeta \frac{\epsilon}{2} K_1 c_T (c_S + 1)) \norm{\partial_1 C^\epsilon}_{L^2}^2  + (K_2 - \zeta \frac{\epsilon}{2} K_1 c_T (c_S + 1)) \norm{\partial_2 C^\epsilon}_{L^2}^2 \\
    & + (K_3 - \zeta \frac{\epsilon}{2} K_1 c_T (c_S + 1)) \norm{\partial_3 C^\epsilon}_{L^2}^2 \bigg] \diff \tau \\
    & \leq \frac{1}{\zeta} K_1 \frac{\epsilon}{2} \int_0 ^ t \norm{\partial_1 C^\epsilon}^2 _{L^2 (\Gamma_G)} + \frac{1}{2} (\norm{\mathbf{u}_H^\epsilon (0)}_{L^2}^2 + \epsilon ^ 2 \norm{u_3^\epsilon (0)}_{L^2}^2 + \norm{C^\epsilon (0)}_{L^2}^2) \\
    & + \int_0^t (S^\epsilon,C^\epsilon) \diff \tau
    \end{split}
  \end{equation*}
\end{frame}

\begin{frame}{Proof steps - Part 1}
  \begin{enumerate}
    \item The beginning step of the proof is to obtain the energy inequality that yields a priori estimates such as $L_t^\infty L_x^2$ boundedness for $C^\epsilon;$ $L_t^2 L_x^2$ boundedness for $\sqrt{\epsilon} \partial_1 C^\epsilon, \partial_2 C^\epsilon, \partial_3 C^\epsilon.$
  
    \item Once we have boundedness, we have the weak convergence of the linear terms, while the weak convergence of the nonlinear terms not including the pollution term had been shown previously by Az\'erad and Guill\'en (SIAM Journal on Mathematical Analysis, 2001).
    
    \item We prove the convergence of the term $(\mathbf{u}^\epsilon C^\epsilon, \nabla \tilde{C})$ by using a compactness result. It can be proved that $u_i^\epsilon C^\epsilon \rightharpoonup  u_i C$ using the following compactness theorem:
  \end{enumerate}
\end{frame}

\begin{frame}{Proof steps - Part 2}

    Compactness theorem.
    
\end{frame}

\section{Visualization - FEM}

\begin{frame}{OCEAN}

    About the ocean BCs and scenario.
    
\end{frame}

\begin{frame}[fragile]
\frametitle{Title of the frame}
\scriptsize

\begin{lstlisting}[frame=single,style=base]
F1_anisotropic =
(1/k)*( (@u1@ - u_prev1)*v1 + (u2 - u_prev2)*v2 + 
eps*eps*(u3 - u_prev3)*v3 ) * dx +
inner(u_prev, grad(u_prev1)) * v1 * dx +
inner(u_prev, grad(u_prev2)) * v2 * dx +
eps*eps*inner(u_prev, grad(u_prev3)) * v3 * dx +
- alpha * u_prev2 * v1 * dx + alpha * u_prev1 * v2 * dx +
inner(grad(u1),grad(v1)) * dx +
inner(grad(u2),grad(v2)) * dx +
eps*eps*inner(grad(u3),grad(v3)) * dx +
eps * beta * u_prev3 * v1 * dx - eps * beta * u_prev1 * v3 * dx
- inner(f, v) * dx +
- inner(theta, v) * ds(1)
\end{lstlisting}
    
\end{frame}

\end{document}